\section{Symbolic identities and interval certificate specifications}
\label{app:symbolic}
\subsection{Exact algebra}
All rational-map identities in Lemma~\ref{lem:rational} are tested by
clearing denominators on their stated nonzero domain and requiring the
numerator to be the zero polynomial. Sparse coefficients are exact
rationals. The checker also embeds the production propagation and
Jacobian formulas in the independent polynomial arithmetic and compares
each entry with direct differentiation. Pointwise numerical agreement
is not the acceptance criterion.

In canonical coordinates the local symbolic certificate supplies the six
fixed-point residuals, the two characteristic polynomials and Jacobian
determinants, the reciprocal-spectrum identities, \(r(t)r(-t)=1\),
the parameter inverse polynomial and \(\sigma'\).
The characteristic sign argument in Lemma~\ref{lem:tails} then establishes
the relevant dimensions; it is not a floating-point eigenvalue count.
The R0 finite checker separately regenerates 35 polynomial identities
and 15 rational comparisons supporting Appendix~\ref{app:carrier}.

For the rational-coordinate proof there are 48 exact identities.
Their roles are the implemented-map and Jacobian comparison, canonical
conjugacy, both inverse identities, reverser involution and reversibility,
contact covariance and determinant, fixed-point and spectral formulas,
the tangent eigenvector, projected contact differentiation and the
conditional scalar derivative cancellation.
One useful propagated phase identity, obtained by differentiating
\eqref{eq:Fhat} and substituting \eqref{eq:contacttangent}, is
\begin{align*}
 z_\zeta={}&-\frac{B_0(y)}{q^2}x_\zeta+
 \left(\frac1{2q}+\frac{(x-1/2)y}{q^2}
       -\frac{2(x-1/2)^2B_0(y)}{q^3}\right)y_\zeta,\\
 Q_\zeta={}&(1/2-y)z_\zeta.
\end{align*}
The independently differentiated formula must agree exactly before its
interval version is used to tighten a propagated derivative enclosure.
The identities alone do not assert completeness of a matching chart,
symmetry of a full match, or a derivative sign.

\subsection{Exact interval domain and analytic majorants}\label{app:intervalspec}
The outer interval \(T^\Box=[t^\Box_-,t^\Box_+]\) is the outward
evaluation of the two radicals \(a(L),a(U)\) by the frozen interval code.
Its exact rational decimal endpoints are
\begin{center}\footnotesize
\(t^\Box_-=\)\texttt{0.878272945180812199461677854479963547717241304597573912661746000456489010475417113206932884}
\par\smallskip
\(t^\Box_+=\)\texttt{0.878272952173734583819480075317750711672962664979312588545615521238490507114264476586355651}
\end{center}
The local complex rectangle is
\(\{t:|\Re t-t_c|\le10^{-5},\,|\Im t|\le10^{-5}\}\), where \(t_c\)
is the recorded real center. Its positive distance from \(T^\Box\)
is \(\Delta_t\), used in the Cauchy derivative error.
All coefficients are enclosed on this complex rectangle, rather than
only at real sample points.

For completeness the majorants used in Certificate~\ref{cert:local}
can be checked from the coefficient boxes as follows. For a phase
radius \(h_0\) and correction allowance \(e_0\), set
\[
 d_i(h_0)=\sum_{n=1}^N|c_{n,i}|h_0^n,\qquad
 X=|t|+d_x(h_0)+e_0,\quad Y=|t|+d_y(h_0)+e_0,
\]
and let \(q_0=\inf|G(t)|-d_q(h_0)-e_0>0\).
Put
\[
 B_M=(1+Y^2)/4,\qquad
 Z=\frac{Y}{2q_0}+\frac{(X+1/2)B_M}{q_0^2}+\frac12.
\]
The coordinate bounds for \(\widehat F_{\sigma(t)}\) are
\[
 Y,\qquad Z,\qquad
 |\sigma(t)|+YZ+(Y+Z)/2+1+B_M/q_0.
\]
For its Jacobian, define
\[
 a_1=\frac{B_M}{q_0^2},\quad
 a_2=\frac1{2q_0}+\frac{(X+1/2)Y}{2q_0^2},\quad
 a_3=\frac{Y}{2q_0^2}+\frac{2(X+1/2)B_M}{q_0^3}.
\]
An upper row-sum bound is
\[
 L_F=\max\left\{1,\ a_1+a_2+a_3,\quad
 (Y+1/2)(a_1+a_2+a_3)+Z+1/2+\frac{Y}{2q_0}
                                      +\frac{B_M}{q_0^2}\right\}.
\]
Use \(h_0=R_0/m,e_0=0\) to bound the defect on \(R_0=0.2\)
by the sum of the propagation and polynomial bounds. Use
\(h_0=r_0/m,e_0=\epsilon_N/m^{N+1}\) for the contraction bound
on \(r_0=0.01\). Complex divisions and square roots reject a
rectangle meeting a forbidden denominator or branch. The finite gates
are \(m>1\), \(q_0>0\), \(\Delta_t>0\), \(\kappa<1/2\), and
\(\delta_N+\kappa\epsilon_N\le\epsilon_N\), with the quantities defined
in the proof of Certificate~\ref{cert:local}. These finite inequalities
are what turn the polynomial into an analytic enclosure.

The recorded corrections and contractions admit the following outward
rational summaries, independently compared with the exact receipts:
\begin{center}
\begin{tabular}{lcc}
Quantity & Degree 36 & Degree 48\\ \hline
Correction bound & \(<3.414\,10^{-47}\)&\(<8.334\,10^{-63}\)\\
Contraction factor & \(<6.891\,10^{-27}\)&\(<4.409\,10^{-36}\)\\
Complex \(|q|\) lower bound & \(>0.2444\)&\(>0.2444\)\\
\(|\mu|\) lower bound & \(>5.8365\)&\(>5.8365\)\\
Parameter margin & \(>9.996\,10^{-6}\)&\(>9.996\,10^{-6}\)
\end{tabular}
\end{center}

\subsection{Cover and scalar receipts}
Each geometry-tree node records a closed phase interval. A nonterminal
node must have two children whose outer endpoints equal the parent's
and whose inner endpoints agree exactly. Both children must be
nondegenerate. Every node must be reachable once from the root;
cycles, missing children and unused nodes are rejected. A geometry
leaf records, for every iterate, the three state intervals, three
phase-derivative intervals and the complementary-\(q\) interval.
Strict endpoint signs certify the predicates of
Certificate~\ref{cert:geometry}.

Every matching-cover leaf covers \(T^\Box\), its source-phase cell,
and \emph{all} possible target phases in the original envelope.
The projected target interval contains every solution of the
monotone inverse equation. A leaf records either the outside-section
label intervals or the projected \(\HH,\KK,\HH_\zeta\) intervals and
associated derivatives. The accepted terminal predicates are precisely
those in Certificate~\ref{cert:cover}. The separate central tree covers
all of \(C_0\), including cells where the main tree used another
predicate. This is why global central monotonicity does not follow
merely from the two retained central leaves.

The scalar receipt encloses the whole implicit phase family inside
\[
 [0.0005824795297982988,0.0005824795369892252].
\]
The exact derivative interval is
\begin{center}\footnotesize
lower endpoint:\par\smallskip
\texttt{3.48002909287387920997544898183182987998806066768305662514348825306622243384968318433976825}
\par\smallskip
upper endpoint:\par\smallskip
\texttt{3.57068402546763126912045926632115988647904165952465340467067063466386183321219918274406572}
\end{center}
It lies strictly in \((3.4800,3.5707)\).
The receipt also encloses \(\sigma'\) by positive endpoints, computes
the quotient derivative in \(s\), and verifies the two global endpoint
residual signs. The root receipt records each input interval, exact
center, residual interval and output interval in \eqref{eq:newton}.
Independent rational arithmetic checks the quotient enclosure, all
13 intersections, their nesting and strict shrinkage, the final
\(10^{-52}\) inflation, the strict signs and exact width.

\subsection{Source-level certificate locator}
Paths in this table are relative to the component roots identified in
Appendix~\ref{app:crosswalk}; they are available offline in the bundled
input archive.
\begin{longtable}{>{\raggedright\arraybackslash}p{0.17\textwidth}>{\raggedright\arraybackslash}p{0.36\textwidth}>{\raggedright\arraybackslash}p{0.37\textwidth}}
Component & Certificate or checker & Mathematical output\\ \hline
R0 & \path{certificates/finite_exact.json};
\path{src/exact_checks.py} & Tail polynomial identities, finite
comparisons, root signs; analytic construction remains Appendix A.\\
R1 & \path{certificate/candidate/};
\path{src/lower.py}, \path{src/upper.py}
& Valid finite strategy and positive exact word identity.\\
Local matching & \path{certificates/symbolic_certificate.json};
\path{src/symbolic_cert.py}
& Fixed points, characteristic factors, inverse parameter formula.\\
Global matching & \path{certificates/symbolic.json};
\path{src/exact.py}
& 48 rational identities, including production-formula agreement.\\
Global matching & \path{certificates/local_manifold.json};
\path{certificates/local_manifold_refined.json};
\path{src/manifold.py}
& Complex analytic-germ and derivative error bounds, degrees 36/48.\\
Global matching & \path{certificates/phase_geometry.json};
\path{certificates/full_cover.json};
\path{certificates/central_monotonicity.json};
\path{src/global_cover.py}
& Complete valid chart; projected inverse; all off-central
exclusions and full central monotonicity.\\
Global matching & \path{certificates/scalar.json};
\path{certificates/root.json}; \path{src/scalar.py}
& Full-domain derivative, existence signs, 13-step strict bracket.\\
Final closure & \path{replay_all.py};
\path{logs/O8_EXACT_CHECKS.json}
& Frozen-interface composition and rational integration arithmetic.\\
\end{longtable}
